\documentclass{article}
\usepackage[utf8]{inputenc}
\usepackage{geometry}
\geometry{a4paper, margin=1in}
\usepackage{minted}

\title{UAV Software Engineering: ROS 2 Architecture}
\author{Advanced Object-Oriented C++ for ROS 2}
\date{}

\begin{document}

\maketitle

\section{The Fallacy of Global State in Avionics}
In embedded systems, relying on global variables (as seen in basic multi-threading examples) is an anti-pattern. In a multi-rotor UAV, you might need to instantiate four identical motor controller nodes. If state is global, these nodes will overwrite each other's memory (Data Races). 

Object-Oriented Programming (OOP) provides \textbf{Encapsulation}. By moving publishers, subscribers, queues, and mutexes into the \texttt{private} scope of a C++ class, each node maintains its own isolated memory footprint. The class constructor safely initializes the DDS entities, and the destructor (following the RAII paradigm) cleanly destroys network endpoints if the node crashes or is shut down.

\section{Inheritance and Shared Ownership (\texttt{rclcpp::Node})}
In ROS 2, every application node strictly inherits from the base class \texttt{rclcpp::Node}. This inheritance provides the derived class with an integrated middleware interface (logging, parameters, time, and network topology).

Crucially, ROS 2 manages the lifecycle of these objects using smart pointers. A ROS 2 node is rarely instantiated on the stack. Instead, it is instantiated on the heap using \texttt{std::make\_shared<MyNode>()}. This is mandatory because the ROS 2 \textbf{Executor} (the engine that spins the node to check for incoming network traffic) requires shared ownership of the node to ensure it is not destroyed while a network callback is still executing.

\section{Hardware Abstraction and the Virtual Table (VTable) Overhead}
OOP allows for Hardware Abstraction Layers (HAL) through \textbf{Polymorphism}. A flight controller might define an abstract base class \texttt{BaseGPS} with a pure virtual function \texttt{virtual void read\_coordinates() = 0;}. Derived classes (e.g., \texttt{UbloxGPS} or \texttt{SeptentrioGPS}) override this function. The core control loop only interacts with \texttt{BaseGPS} pointers, making the drone hardware-agnostic.

However, in hard real-time C++, polymorphism introduces overhead. When a \texttt{virtual} function is called, the compiler cannot resolve the address at compile-time (Dynamic Dispatch). Instead, it looks up the function address in a hidden array of pointers called the \textbf{Virtual Table (vtable)} via a Virtual Pointer (vptr) injected into the object's memory. In ultra-low-latency loops (e.g., 4kHz IMU reads), navigating the vtable can cause CPU cache misses. Modern UAV software mitigates this by using the \texttt{final} keyword to explicitly terminate inheritance trees, allowing the compiler to devirtualize calls.

\section{Event-Driven Execution: \texttt{std::bind} vs Lambdas}
ROS 2 is event-driven: the middleware triggers a callback when a DDS packet arrives. Pointing the middleware to a class member function is complex because member functions require a hidden \texttt{this} pointer to identify the specific object instance.

Historically, ROS 2 used \texttt{std::bind} to attach the \texttt{this} pointer and the message placeholders (\texttt{std::placeholders::\_1}) to the callback. However, modern C++ (C++14/17) strongly prefers \textbf{Lambda Expressions} capturing \texttt{[this]}. Lambdas are generally superior in robotics because:
\begin{enumerate}
    \item They are easier for the compiler to inline, reducing function call overhead.
    \item They provide strict compile-time type checking, whereas \texttt{std::bind} errors often manifest as cryptic template errors.
    \item They allow for inline data preprocessing (e.g., converting GPS coordinates from degrees to radians) before passing the data to the main class logic.
\end{enumerate}

\section{The Pimpl Idiom (Pointer to Implementation)}
A critical architectural pattern in professional UAV software is the \textbf{Pimpl Idiom}. When building a ROS 2 node that interacts with proprietary hardware drivers or complex third-party libraries (like OpenCV or a specific flight controller SDK), exposing these dependencies in the node's header file (\texttt{.hpp}) forces a recompilation of the entire software stack whenever a minor detail changes.

By forward-declaring a private implementation struct (e.g., \texttt{struct Impl;}) in the header and holding a \texttt{std::unique\_ptr<Impl>} in the class, the true implementation details are hidden inside the \texttt{.cpp} file. This guarantees stable Application Binary Interfaces (ABI) and drastically reduces compile times in massive drone codebases.

\section{Resource Management: The Rule of Delete}
A UAV software node often encapsulates exclusive system resources: DDS network endpoints, active \texttt{std::thread} instances, or hardware serial ports (UART/I2C). These resources are strictly non-copyable. If a C++ class managing a thread is passed by value, the compiler implicitly generates a copy constructor, cloning the state and leading to catastrophic hardware collisions or double-free memory corruption during destruction.

To enforce compile-time safety, professional C++ avionics code explicitly deletes the copy constructor and copy assignment operator (a subset of the Rule of Five). Any attempt to accidentally copy the node will immediately halt the compilation:
\begin{minted}{cpp}
class TelemetryNode {
public:
    // Delete copy semantics
    TelemetryNode(const TelemetryNode&) = delete;
    TelemetryNode& operator=(const TelemetryNode&) = delete;
};
\end{minted}

\section{Zero-Cost Abstraction: The CRTP Pattern}
While Section 3 highlights the latency overhead of dynamic polymorphism (vtable lookups)[cite: 3], 
embedded software still requires Hardware Abstraction Layers. The industry solution for zero-cost abstraction 
is \textbf{Static Polymorphism} via the \textbf{Curiously Recurring Template Pattern (CRTP)}.

Instead of using \texttt{virtual} keywords, the base class is implemented as a template, and the derived class passes itself as 
the template argument:

\begin{minted}{cpp}
template <typename Derived>
class BaseSensor {
public:
    void read_data() {
        // Compile-time static cast, zero runtime overhead
        static_cast<Derived*>(this)->impl_read_data(); 
    }
};

class FastIMU : public BaseSensor<FastIMU> {
public:
    void impl_read_data() { /* hardware specific logic */ }
};
\end{minted}

This paradigm provides the structural benefits of OOP interfaces while allowing the compiler to aggressively inline the function calls, 
making it the gold standard for high-frequency control loops (e.g., 400Hz+ motor mixing).

\end{document}